\section{Second Pass: Functions as Maps of Structure}
\chapterline

In Chapter~31, we treated sets as static containers. But mathematical spaces usually carry additional structure: operations that combine elements, metrics that measure distance, or symmetries that preserve shapes. Functions between such spaces are not merely input-output rules; they are maps that preserve or reflect that structure.

In elementary algebra, a function is often introduced as a calculation formula, such as $f(x) = x^2 + 1$. In higher mathematics, this shifts to a structural perspective. A function between structured spaces is viewed as a \emph{morphism}---a structure-preserving map between two spaces. The central question is: \emph{what underlying geometric or algebraic relationships are preserved?}

\begin{historybox}
In 1945, two mathematicians working at the intersection of algebra and geometry---Saunders Mac Lane and Samuel Eilenberg---published a landmark paper titled \emph{General Theory of Natural Equivalences}. They were confronting a concrete obstacle: in algebraic topology (the field that analyzes geometric shapes by translating them into algebraic structures), mathematicians calculated homology groups (algebraic invariants that count independent holes, tunnels, and voids in a space). Different topological constructions produced isomorphic groups, but every isomorphism depended on choosing arbitrary coordinate bases. They had no rigorous language to state that two geometric constructions were ``naturally'' identical without making an arbitrary choice.

To solve this, Eilenberg and Mac Lane introduced a new perspective. Instead of inspecting mathematical objects solely by their elements, they formulated \emph{Category Theory}, representing mathematical structures through collections of \emph{objects} connected by directed \emph{arrows} (morphisms). Mac Lane observed that categories were introduced to define functors (maps that translate categories into categories while preserving composition), and functors were introduced to define natural transformations (systematic mappings between functors).

When category theory was first presented, some algebraists referred to it as ``generalized abstract nonsense,'' viewing arrow diagrams as purely formal bookkeeping. The perspective proved fruitful in the 1950s and 1960s when Alexander Grothendieck placed algebraic geometry on categorical foundations, resolving conjectures by characterizing objects through the morphisms probing them. Today, categorical thinking is widely used across algebraic topology, mathematical physics, and programming language semantics. An active area of research is \emph{higher category theory} (such as $(\infty,1)$-categories and homotopy type theory), where arrows connect objects, higher arrows connect arrows, and the hierarchy continues.
\end{historybox}

\subsection{Preserving Structure Across Different Mathematical Worlds}

Depending on the discipline, sets come equipped with additional features---distances, symmetries, algebraic operations, or notions of openness. A structured function is one that respects those features:

\begin{itemize}
  \item \textbf{Analysis (Metric Spaces)}: A metric space is a set equipped with a distance function $d$. A function $f: X \to Y$ is continuous if it respects nearness. If a sequence of points $(x_n)$ converges to $x$, then their images $f(x_n)$ converge to $f(x)$. Nearby inputs produce nearby outputs.
  \item \textbf{Abstract Algebra (Groups and Rings)}: A group is a set equipped with an associative operation, an identity element, and inverses. A group homomorphism $\phi: G \to H$ respects operations:
  \[
  \phi(a \cdot_G b) = \phi(a) \cdot_H \phi(b).
  \]
  Multiplying inside $G$ and then mapping produces the exact same result as mapping first and multiplying inside $H$.
  \item \textbf{Topology}: A topological space generalizes distance by specifying which subsets count as ``open neighborhoods.'' A homeomorphism $f: X \to Y$ is a continuous bijection with a continuous inverse. It deforms a space without tearing or gluing, keeping connectivity intact.
  \item \textbf{Differential Geometry}: A manifold is a space that locally resembles flat Euclidean space $\R^n$. A smooth map between manifolds preserves the ability to do calculus; directional derivatives and tangent spaces transform smoothly.
\end{itemize}

\begin{intuitionbox}
To see how a homomorphism preserves operations in everyday math, picture the exponential function:
\[
\exp: (\R, +) \longrightarrow (\R_{>0}, \times), \qquad \exp(x) = e^x.
\]
The source world $(\R, +)$ is the real line where elements combine by addition. The target world $(\R_{>0}, \times)$ is the set of positive numbers combining by multiplication.

Pick two numbers in the source, say $x = 2$ and $y = 3$. You have two ways to reach the target:
\begin{enumerate}
  \item \textbf{Add first, then map}: In the source, $2 + 3 = 5$. Evaluating $\exp$ gives $\exp(5) = e^5$.
  \item \textbf{Map first, then multiply}: Evaluating $\exp$ on each term first gives $\exp(2) = e^2$ and $\exp(3) = e^3$. Multiplying them in the target gives $e^2 \cdot e^3 = e^5$.
\end{enumerate}
Because $e^{x+y} = e^x \cdot e^y$, both routes end at the exact same value. Exponentiation translates the addition of the real line into the multiplication of positive numbers while preserving the algebraic structure.
\end{intuitionbox}

Figure~\ref{32:fig-homomorphism-commute} illustrates this two-path principle: a homomorphism guarantees that combining in the source and then mapping yields the exact same destination as mapping first and combining in the target.

\begin{figure}[htbp]
\centering
\begin{tikzpicture}[scale=0.95, >=Latex]
  % Left Panel: Abstract Homomorphism
  \begin{scope}[xshift=0cm]
    \node[anchor=south, font=\small\bfseries] at (2.7, 2.3) {Abstract Homomorphism $\phi: G \to H$};
    \draw[rounded corners=4pt, fill=gray!3, draw=gray!25] (-0.2, -1.8) rectangle (5.6, 2.3);
    
    \node[font=\footnotesize\bfseries, text=sourceblue!90!black] (tl) at (0.8, 1.3) {$(a, b)$};
    \node[font=\footnotesize\bfseries, text=sourceblue!90!black] (tr) at (4.6, 1.3) {$a \cdot_G b$};
    \node[font=\footnotesize\bfseries, text=targetgreen!90!black] (bl) at (0.8, -0.7) {$(\phi(a), \phi(b))$};
    \node[font=\footnotesize\bfseries, text=targetgreen!90!black] (br) at (4.6, -0.7) {$\phi(a) \cdot_H \phi(b)$};
    
    \draw[->, thick, draw=sourceblue!80!black] (tl) -- node[above, font=\scriptsize] {combine in $G$} (tr);
    \draw[->, thick, draw=structurepurple!80!black] (tl) -- node[left, font=\scriptsize] {$(\phi, \phi)$} (bl);
    \draw[->, thick, draw=structurepurple!80!black] (tr) -- node[right, font=\scriptsize] {$\phi$} (br);
    \draw[->, thick, draw=targetgreen!80!black] (bl) -- node[below, font=\scriptsize] {combine in $H$} (br);
    
    \node at (2.7, 0.3) [font=\large, text=gray!60] {$\circlearrowright$};
    \node at (2.7, -0.15) [font=\tiny\bfseries, text=gray!70!black] {both paths agree};
    \node[font=\scriptsize, align=center, text=gray!80!black] at (2.7, -1.45) {$\phi(a \cdot_G b) = \phi(a) \cdot_H \phi(b)$};
  \end{scope}

  % Right Panel: Concrete Exponentiation
  \begin{scope}[xshift=6.8cm]
    \node[anchor=south, font=\small\bfseries] at (2.7, 2.3) {Concrete Model: $(\R, +) \to (\R_{>0}, \times)$};
    \draw[rounded corners=4pt, fill=gray!3, draw=gray!25] (-0.2, -1.8) rectangle (5.6, 2.3);
    
    \node[font=\footnotesize\bfseries, text=sourceblue!90!black] (etl) at (0.8, 1.3) {$(2, 3)$};
    \node[font=\footnotesize\bfseries, text=sourceblue!90!black] (etr) at (4.6, 1.3) {$2 + 3 = 5$};
    \node[font=\footnotesize\bfseries, text=targetgreen!90!black] (ebl) at (0.8, -0.7) {$(e^2, e^3)$};
    \node[font=\footnotesize\bfseries, text=targetgreen!90!black] (ebr) at (4.6, -0.7) {$e^2 \cdot e^3 = e^5$};
    
    \draw[->, thick, draw=sourceblue!80!black] (etl) -- node[above, font=\scriptsize] {add} (etr);
    \draw[->, thick, draw=structurepurple!80!black] (etl) -- node[left, font=\scriptsize] {$\exp$} (ebl);
    \draw[->, thick, draw=structurepurple!80!black] (etr) -- node[right, font=\scriptsize] {$\exp$} (ebr);
    \draw[->, thick, draw=targetgreen!80!black] (ebl) -- node[below, font=\scriptsize] {multiply} (ebr);
    
    \node at (2.7, 0.3) [font=\large, text=gray!60] {$\circlearrowright$};
    \node at (2.7, -0.15) [font=\tiny\bfseries, text=gray!70!black] {commutes};
    \node[font=\scriptsize, align=center, text=gray!80!black] at (2.7, -1.45) {$\exp(2 + 3) = \exp(2) \cdot \exp(3)$};
  \end{scope}
\end{tikzpicture}
\caption{The commutative geometry of operation preservation. A homomorphism $\phi$ guarantees that combining in the domain and then translating yields the exact same outcome as translating first and combining in the codomain. On the right, the exponential map translates real addition into positive multiplication without distortion.}
\label{32:fig-homomorphism-commute}
\end{figure}

\begin{recoverybox}
\textbf{Pause and think}: Consider the shift map $s: \R \to \R$ given by $s(x) = x + 1$. Is $s$ a group homomorphism from $(\R, +)$ to $(\R, +)$?

Test two inputs, say $x = 1$ and $y = 2$:
\begin{itemize}
  \item Adding first gives $s(1 + 2) = s(3) = 4$.
  \item Shifting first and adding gives $s(1) + s(2) = 2 + 3 = 5$.
\end{itemize}
They disagree ($4 \neq 5$)! Shifting by $1$ moves points, but it disrupts addition because it fails to keep the identity anchored: $s(0) = 1 \neq 0$. A true group homomorphism must always carry the identity to the identity.
\end{recoverybox}

Mathematicians frequently understand an object not by staring at its individual elements in isolation, but by observing how other objects map into it, and how it maps into other objects.

\subsection{Subtle Traps: Basis-Dependent Maps and One-Way Continuity}

Thinking in terms of structure-preserving maps clears up two traps that trip up almost everyone at first:

\begin{itemize}
  \item \textbf{The Pitfall of Non-Canonical Constructions}: A common difficulty is attempting to define a mathematical relationship by choosing arbitrary coordinates. For instance, in linear algebra, the dual space $V^*$ is the set of linear functionals (maps that take a vector and output a scalar). Every finite-dimensional vector space $V$ has the same dimension as $V^*$ and is therefore isomorphic to it. But any such isomorphism requires choosing a basis; if you pick a different basis, the isomorphism changes. In contrast, each vector $v \in V$ acts naturally as an evaluation functional on $V^*$: $\operatorname{eval}_v(\varphi) = \varphi(v)$. The resulting map $v \mapsto \operatorname{eval}_v$ from $V$ to its double dual $V^{**}$ requires no choice of basis---it is \emph{canonical}. Morphism thinking prioritizes coordinate-independent mappings that reflect intrinsic structure.
  \item \textbf{Continuous Bijections Need Not Be Isomorphisms}: In set theory, a bijection is an invertible map. But in geometry and topology, a continuous map can be strictly bijective without being an isomorphism (homeomorphism). For example, the map $f: [0, 1) \to S^1$ given by $f(t) = e^{2\pi i t}$ wraps the half-open interval onto the unit circle in the complex plane. Every point on the circle has exactly one preimage, and the forward map is continuous. But its inverse is discontinuous at $(1, 0)$: points on the circle just below the positive horizontal axis have preimages near $t = 1$, while $(1, 0)$ itself has preimage $t = 0$. The inverse fails to preserve open sets. A true structural isomorphism must have an inverse that also preserves the structure of the category.
\end{itemize}

Figure~\ref{32:fig-wrapping-tearing} reveals why bijectivity is insufficient: the inverse map maps nearby points on the circle to points separated by unit distance on the interval, failing to preserve the topology.

\begin{figure}[htbp]
\centering
\begin{tikzpicture}[scale=0.95, >=Latex]
  % Left: Half-open interval [0, 1)
  \begin{scope}[xshift=0cm]
    \node[font=\small\bfseries, text=sourceblue!90!black] at (1.8, 2.3) {1. Domain: Interval $[0, 1)$};
    \draw[rounded corners=4pt, fill=gray!3, draw=gray!25] (-0.3, -2.1) rectangle (3.9, 2.1);
    
    \draw[thick, gray!50, ->] (0.0, 0.3) -- (3.6, 0.3);
    \draw[line width=2.5pt, sourceblue] (0.4, 0.3) -- (3.0, 0.3);
    
    \fill[sourceblue] (0.4, 0.3) circle (2.5pt);
    \node[below=3pt, font=\scriptsize\bfseries, text=sourceblue!90!black] at (0.4, 0.3) {$t = 0$};
    
    \draw[thick, sourceblue, fill=white] (3.1, 0.3) circle (2.5pt);
    \node[below=3pt, font=\scriptsize, text=gray!70!black] at (3.1, 0.3) {$1$ (open)};
    
    \fill[targetgreen] (2.7, 0.3) circle (2.2pt);
    \node[above=3pt, font=\tiny\bfseries, text=targetgreen!90!black] at (2.7, 0.3) {$t = 0.99$};
    
    \draw[<->, thick, trapred] (0.4, 0.9) -- (2.7, 0.9)
      node[midway, above, font=\tiny\bfseries, text=trapred] {distance $\approx 1$ (far apart)};
    
    \node[font=\scriptsize, align=center, text=gray!80!black] at (1.8, -1.2) {Points $0$ and $0.99$ are\\at opposite ends of interval};
  \end{scope}

  % Middle transition arrows
  \begin{scope}[xshift=4.1cm]
    \draw[->, thick, draw=structurepurple!90!black] (0.1, 0.6) to[bend left=15]
      node[above, font=\scriptsize\bfseries, text=structurepurple!90!black] {$f(t) = e^{2\pi i t}$} (2.3, 0.6);
    
    \draw[<-, thick, draw=trapred!90!black, dashed] (0.1, -0.4) to[bend right=15]
      node[below=2pt, font=\scriptsize\bfseries, text=trapred!90!black] {$f^{-1}$ fails continuity} (2.3, -0.4);
  \end{scope}

  % Right: Unit circle S^1
  \begin{scope}[xshift=6.7cm]
    \node[font=\small\bfseries, text=targetgreen!90!black] at (2.0, 2.3) {2. Codomain: Circle $S^1 \subset \C$};
    \draw[rounded corners=4pt, fill=gray!3, draw=gray!25] (-0.3, -2.1) rectangle (4.3, 2.1);
    
    \draw[gray!35] (0.6, 0.4) -- (3.4, 0.4);
    \draw[gray!35] (2.0, -1.0) -- (2.0, 1.8);
    
    \draw[line width=2pt, sourceblue!80!black] (2.0, 0.4) circle (0.95);
    
    \fill[sourceblue] (2.95, 0.4) circle (2.5pt);
    \node[right=2pt, font=\tiny\bfseries, text=sourceblue!90!black] at (2.95, 0.55) {$f(0) = 1$};
    
    \fill[targetgreen] ({2.0 + 0.95*cos(-18)}, {0.4 + 0.95*sin(-18)}) circle (2.2pt);
    \node[right=2pt, font=\tiny\bfseries, text=targetgreen!90!black] at ({2.0 + 0.95*cos(-18)}, {0.4 + 0.95*sin(-18) - 0.12}) {$f(0.99)$};
    
    \node[font=\scriptsize, align=center, text=gray!80!black] at (2.0, -1.4) {Images $f(0)$ and $f(0.99)$ are close;\\inverse values remain far apart};
  \end{scope}
\end{tikzpicture}
\caption{Why a continuous bijection is not necessarily an isomorphism (homeomorphism). The wrapping map $f(t) = e^{2\pi i t}$ sends $[0, 1)$ bijectively onto the unit circle $S^1$. Although $f$ is continuous, its inverse $f^{-1}$ is discontinuous at $(1, 0)$: points arbitrarily close on the circle just above and below the positive real axis map to opposite ends of the interval $[0, 1)$.}
\label{32:fig-wrapping-tearing}
\end{figure}

\subsection{Commutative Diagrams: The Language of Compatibility}

When several maps interact, their relationships are best visualized using \emph{commutative diagrams}. A commutative diagram is a diagram of mathematical objects and arrows where any two directed paths between the same pair of objects yield the same composite function.

The definition of a group homomorphism $\phi: G \to H$ can be expressed as a commutative diagram. In the source group $G$, the operation is a function $\mu_G: G \times G \to G$ given by $\mu_G(a, b) = a \cdot_G b$. In the target group $H$, the operation is $\mu_H: H \times H \to H$. Saying that $\phi(a \cdot_G b) = \phi(a) \cdot_H \phi(b)$ means that the following diagram commutes:
\[
\begin{array}{ccc}
G \times G & \xrightarrow{\quad \mu_G \quad} & G \\
\Big\downarrow{\scriptstyle (\phi, \phi)} & & \Big\downarrow{\scriptstyle \phi} \\
H \times H & \xrightarrow{\quad \mu_H \quad} & H
\end{array}
\]
Starting at $(a, b)$ in the top-left, going across and down produces $\phi(a \cdot_G b)$. Going down and across produces $\phi(a) \cdot_H \phi(b)$. Commutativity is the visual syntax of algebraic compatibility.

\paragraph{The Canonical Decomposition of a Function.}
Consider an arbitrary function $f: A \to B$. If we define an equivalence relation where two inputs are related whenever they produce the same output ($x \sim y \iff f(x) = f(y)$), the domain $A$ partitions into fibers $A/\sim$. Every function factors cleanly into three separate stages, illustrated in Figure~\ref{32:fig-canonical-factorization}:
\begin{enumerate}
  \item \textbf{Projection ($\pi$)}: The canonical map $\pi: A \to A/\sim$ defined by $\pi(x) = [x]$. This map is surjective; it sends each element to its equivalence class.
  \item \textbf{Induced bijection ($\bar{f}$)}: The map $\bar{f}: A/\sim \to \operatorname{im}(f)$ defined by $\bar{f}([x]) = f(x)$. This map is both injective and surjective; it is a bijection between fibers and attained values.
  \item \textbf{Inclusion ($\iota$)}: The map $\iota: \operatorname{im}(f) \to B$ defined by $\iota(y) = y$. This map is injective; it places the image inside the codomain $B$.
\end{enumerate}

\begin{figure}[htbp]
\centering
\begin{tikzpicture}[scale=1.15, >=Stealth]
  \tikzset{
    objnode/.style={draw=sourceblue!80!black, fill=sourceblue!10, very thick, rounded corners=3pt, inner sep=6pt, font=\normalsize\bfseries},
    quotnode/.style={draw=quotientorange!80!black, fill=quotientorange!12, very thick, rounded corners=3pt, inner sep=6pt, font=\normalsize\bfseries},
    tarnode/.style={draw=targetgreen!80!black, fill=targetgreen!12, very thick, rounded corners=3pt, inner sep=6pt, font=\normalsize\bfseries},
    maparrow/.style={->, thick, shorten >=2pt, shorten <=2pt},
    annot/.style={font=\footnotesize, text=gray!80!black}
  }

  % Nodes
  \node[objnode] (A) at (0, 3) {$A$};
  \node[tarnode] (B) at (6, 3) {$B$};
  \node[quotnode] (Q) at (0, 0) {$A/\sim$};
  \node[tarnode] (im) at (6, 0) {$\operatorname{im}(f)$};

  % Arrows
  \draw[maparrow, draw=sourceblue!80!black] (A) -- node[above, font=\small] {$f$ (original map)} (B);
  \draw[maparrow, draw=structurepurple!80!black] (A) -- node[left, font=\small] {$\pi$ (projection)} (Q);
  \draw[maparrow, draw=quotientorange!90!black, dashed] (Q) -- node[below, font=\small] {$\bar{f}$ (induced bijection)} (im);
  \draw[maparrow, draw=targetgreen!80!black] (im) -- node[right, font=\small] {$\iota$ (inclusion)} (B);

  % Commutativity symbol in center
  \node at (3, 1.6) [font=\large, text=gray!60] {$\circlearrowright$};
  \node at (3, 1.1) [font=\footnotesize, text=gray!70!black] {$f = \iota \circ \bar{f} \circ \pi$};

  % Side annotations with generous clearance
  \node[annot, left=0.35cm of Q, align=right] {Fibers mapped to\\single elements};
  \node[annot, right=0.35cm of im, align=left] {Exact set of\\attained values};
\end{tikzpicture}
\caption{The canonical factorization of a function $f: A \to B$. Any mapping can be factored into three canonical stages: projecting each element to its fiber ($\pi$), mapping each fiber to its image via a bijection ($\bar{f}$), and viewing the image as a subset of the codomain via inclusion ($\iota$). Both paths yield the same function: $f = \iota \circ \bar{f} \circ \pi$.}
\label{32:fig-canonical-factorization}
\end{figure}

In this diagram, taking the top route ($f$) produces the exact same outcome as taking the perimeter detour: project down to the quotient via $\pi$, apply the induced bijection $\bar{f}$, and include the image into $B$ via $\iota$. The equation
\[
f = \iota \circ \bar{f} \circ \pi
\]
expresses commutativity algebraically. Every function factors into these three structural stages: a projection to the quotient, an induced bijection, and an inclusion.

\subsection{Kernels: Measuring What a Map Forgets}

Every non-injective map collapses some information. In linear algebra, a matrix transformation may map a three-dimensional space onto a flat plane; the subspace mapped to the zero vector is the \emph{null space}.

In group theory, where $G$ and $H$ are groups with identity elements $e_G$ and $e_H$, the \emph{kernel} of a homomorphism $\phi: G \to H$ is defined as:
\[
\ker(\phi) = \{g \in G : \phi(g) = e_H\}.
\]
The kernel is the algebraic analogue of the null space: it collects every element of $G$ that maps to the identity in $H$.

\begin{intuitionbox}
The size of $\ker(\phi)$ measures how much information $\phi$ collapses. If $\ker(\phi) = \{e_G\}$, no two distinct elements map to the same image, and $\phi$ is injective. If $\ker(\phi)$ contains non-identity elements, each fiber contains multiple elements.
\end{intuitionbox}

\begin{theorem}\label{32:thm-kernel-injective}
A group homomorphism $\phi: G \to H$ is injective if and only if $\ker(\phi) = \{e_G\}$.
\end{theorem}

\begin{proof}
First, observe that every group homomorphism preserves the identity:
\[
\phi(e_G) = \phi(e_G \cdot e_G) = \phi(e_G) \cdot \phi(e_G).
\]
Multiplying both sides by $\phi(e_G)^{-1}$ in $H$ shows that $\phi(e_G) = e_H$.

$(\implies)$ Suppose $\phi$ is injective. If $g \in \ker(\phi)$, then $\phi(g) = e_H = \phi(e_G)$. Because $\phi$ is injective, this immediately forces $g = e_G$. Hence $\ker(\phi) = \{e_G\}$.

$(\impliedby)$ Suppose $\ker(\phi) = \{e_G\}$. Let $a, b \in G$ satisfy $\phi(a) = \phi(b)$. Multiply both sides on the right by $\phi(b)^{-1}$ in $H$:
\[
\phi(a)\phi(b)^{-1} = e_H.
\]
Because homomorphisms preserve inverses ($\phi(b)\phi(b^{-1}) = \phi(b b^{-1}) = \phi(e_G) = e_H \implies \phi(b)^{-1} = \phi(b^{-1})$) and products, we can rewrite this as:
\[
\phi(a)\phi(b^{-1}) = \phi(ab^{-1}) = e_H.
\]
This implies $ab^{-1} \in \ker(\phi)$. Since the kernel contains only the identity $e_G$, we have $ab^{-1} = e_G$. Multiplying by $b$ on the right gives $a = b$. Thus $\phi$ is injective.
\end{proof}

\begin{recoverybox}
\textbf{Pause and think}: Why does testing injectivity in group theory require checking only one element (the identity), whereas in general set theory you must check every pair of elements?

Because groups are homogeneous: every point looks algebraically identical to every other point. If two distinct inputs collide ($\phi(a) = \phi(b)$), we can always slide the collision back to the identity by computing $ab^{-1}$: $\phi(ab^{-1}) = \phi(a)\phi(b)^{-1} = e_H$. If the identity has no collisions, no point anywhere in the group has collisions.
\end{recoverybox}

\subsection{Quotient Groups and the First Isomorphism Theorem}

Once we identify the kernel $\ker \phi$, we can construct the quotient group $G/\ker\phi$, identifying elements that have the same image under $\phi$.

Why does the kernel dictate the shape of every other fiber? Suppose two elements $x, a \in G$ produce the exact same output: $\phi(x) = \phi(a)$. Then:
\[
\phi(a)^{-1}\phi(x) = e_H \implies \phi(a^{-1}x) = e_H \implies a^{-1}x \in \ker(\phi).
\]
Setting $k = a^{-1}x \in \ker(\phi)$, we obtain $x = ak$. Conversely, for any $k \in \ker(\phi)$, the homomorphism property ensures:
\[
\phi(ak) = \phi(a)\phi(k) = \phi(a)e_H = \phi(a).
\]
Therefore, the full preimage of the point $\phi(a)$ is precisely the left coset:
\[
a\ker(\phi) = \{ak : k \in \ker(\phi)\}.
\]
Every single fiber of $\phi$ is an exact translated copy of the kernel. If $\ker(\phi)$ contains $k$ elements, every fiber across the entire group contains exactly $k$ elements. The domain $G$ is sliced into parallel coset sheets, as pictured in Figure~\ref{32:fig-first-isomorphism}.

\begin{figure}[htbp]
\centering
\begin{tikzpicture}[scale=1.0, >=Stealth]
  \tikzset{
    domainbox/.style={draw=sourceblue!80!black, fill=sourceblue!5, very thick, rounded corners=6pt},
    fiberbox/.style={draw=sourceblue!70!black, fill=sourceblue!15, thick, rounded corners=3pt},
    kerbox/.style={draw=quotientorange!90!black, fill=quotientorange!20, very thick, rounded corners=3pt},
    targbox/.style={draw=targetgreen!80!black, fill=targetgreen!5, very thick, rounded corners=6pt},
    imbox/.style={draw=targetgreen!70!black, fill=targetgreen!15, thick, rounded corners=4pt},
    pointdot/.style={circle, fill=black, inner sep=1.8pt},
    quotpoint/.style={circle, fill=quotientorange!90!black, inner sep=2.2pt},
    maparrow/.style={->, thick, shorten >=3pt, shorten <=3pt},
    collapsearrow/.style={->, dashed, thick, shorten >=3pt, shorten <=3pt}
  }

  % Domain G
  \draw[domainbox] (0, 0) rectangle (3.2, 4.4);
  \node[font=\bfseries, text=sourceblue!90!black] at (1.6, 4.05) {Group $G$};

  % Fiber 1: a ker(phi)
  \draw[fiberbox] (0.25, 2.7) rectangle (2.95, 3.6);
  \node[pointdot, label={above:\footnotesize $a$}] at (1.0, 2.95) {};
  \node[font=\footnotesize, text=sourceblue!90!black] at (2.2, 3.15) {$a \ker\phi$};

  % Fiber 2: Kernel ker(phi)
  \draw[kerbox] (0.25, 1.5) rectangle (2.95, 2.4);
  \node[pointdot, label={above:\footnotesize $e_G$}] at (1.0, 1.75) {};
  \node[font=\footnotesize\bfseries, text=quotientorange!95!black] at (2.2, 1.95) {$\ker\phi$};

  % Fiber 3: b ker(phi)
  \draw[fiberbox] (0.25, 0.3) rectangle (2.95, 1.2);
  \node[pointdot, label={above:\footnotesize $b$}] at (1.0, 0.55) {};
  \node[font=\footnotesize, text=sourceblue!90!black] at (2.2, 0.75) {$b \ker\phi$};

  % Target H
  \draw[targbox] (7.2, 0) rectangle (10.4, 4.4);
  \node[font=\bfseries, text=targetgreen!90!black] at (8.8, 4.05) {Group $H$};

  % Image im(phi)
  \draw[imbox] (7.45, 0.3) rectangle (10.15, 3.6);
  \node[font=\scriptsize\bfseries, text=targetgreen!80!black] at (8.8, 0.55) {$\operatorname{im}(\phi)$};

  % Target points
  \node[pointdot, label={right:\footnotesize $\phi(a)$}] (pha) at (8.3, 3.15) {};
  \node[pointdot, label={right:\footnotesize $e_H = \phi(e_G)$}] (phe) at (8.3, 1.95) {};
  \node[pointdot, label={right:\footnotesize $\phi(b)$}] (phb) at (8.3, 1.15) {};

  % Collapse arrows from fibers to points
  \draw[collapsearrow, draw=sourceblue!70!black] (2.95, 3.15) -- (pha);
  \draw[collapsearrow, draw=quotientorange!90!black] (2.95, 1.95) -- (phe);
  \draw[collapsearrow, draw=sourceblue!70!black] (2.95, 0.75) -- (phb);

  % Top morphism arrow
  \draw[maparrow, draw=structurepurple!90!black] (3.4, 4.05) -- node[above, font=\small] {$\phi$ (homomorphism)} (7.0, 4.05);

  % Annotation above dashed arrows
  \node[font=\scriptsize, text=gray!70!black, align=center] at (5.2, 3.65) {coset fibers collapse\\to single points};

  % Bottom Quotient representation
  \node[draw=quotientorange!80!black, fill=quotientorange!8, very thick, rounded corners=4pt, inner sep=6pt] (Qbox) at (5.2, -1.6) {
    Quotient Group $G/\ker\phi = \{a \ker\phi, \, \ker\phi, \, b \ker\phi, \, \dots\}$
  };

  % Arrows to and from quotient
  \draw[maparrow, draw=structurepurple!80!black] (1.6, 0) -- node[left, font=\footnotesize] {$\pi$ (projection)} (1.6, -1.25);
  \draw[maparrow, draw=targetgreen!90!black] (8.8, -1.25) -- node[right=4pt, font=\footnotesize] {$\bar{\phi}$ (isomorphism)} (8.8, 0);
\end{tikzpicture}
\caption{The geometry of the First Isomorphism Theorem: $G/\ker(\phi) \cong \operatorname{im}(\phi)$. A homomorphism $\phi$ slices the domain $G$ into parallel coset fibers. The kernel $\ker(\phi)$ is the distinguished fiber that collapses to the identity $e_H$. Every other fiber $g \ker(\phi)$ collapses entirely to the single point $\phi(g)$ in the image. Passing to the quotient $G/\ker(\phi)$ packages each entire fiber into a single point, producing an exact copy of the image.}
\label{32:fig-first-isomorphism}
\end{figure}

\begin{theorem}[The First Isomorphism Theorem for Groups]\label{32:thm-first-isomorphism}
Let $\phi: G \to H$ be a group homomorphism. Then the map
\[
\bar{\phi}: G/\ker(\phi) \longrightarrow \operatorname{im}(\phi), \qquad \bar{\phi}(g\ker\phi) = \phi(g),
\]
is a well-defined group isomorphism. Consequently,
\[
G/\ker(\phi) \cong \operatorname{im}(\phi).
\]
\end{theorem}

\begin{proof}
We check the four required properties in turn:
\begin{enumerate}
  \item \textbf{Well-defined}: Suppose $g\ker\phi = g'\ker\phi$. Then $g' = gk$ for some $k \in \ker\phi$. Evaluating $\bar{\phi}$ gives:
  \[
  \phi(g') = \phi(gk) = \phi(g)\phi(k) = \phi(g)e_H = \phi(g).
  \]
  The value does not depend on the choice of representative.
  \item \textbf{Homomorphism}: Coset multiplication in $G/\ker\phi$ is defined by $(a\ker\phi)(b\ker\phi) = ab\ker\phi$. Thus:
  \[
  \bar{\phi}((a\ker\phi)(b\ker\phi)) = \bar{\phi}(ab\ker\phi) = \phi(ab) = \phi(a)\phi(b) = \bar{\phi}(a\ker\phi)\bar{\phi}(b\ker\phi).
  \]
  \item \textbf{Injective}: Suppose $\bar{\phi}(g\ker\phi) = e_H$. Then $\phi(g) = e_H$, so $g \in \ker\phi$. But this means $g\ker\phi = \ker\phi$, which is the identity element of $G/\ker\phi$. By Theorem~\ref{32:thm-kernel-injective}, $\bar{\phi}$ is injective.
  \item \textbf{Surjective}: By definition, every element of $\operatorname{im}(\phi)$ has the form $\phi(g)$ for some $g \in G$. Since $\bar{\phi}(g\ker\phi) = \phi(g)$, every image point is reached.
\end{enumerate}
Because $\bar{\phi}$ is a bijective homomorphism, it is an isomorphism.
\end{proof}

\begin{trapbox}
\textbf{The Well-Definedness Trap}: Whenever you define a function whose input is an equivalence class or a coset, your formula inevitably picks a representative: $\bar{\phi}([g]) = \phi(g)$. But $[g]$ contains many elements! If someone else selects $g' \in [g]$, will your formula return the same answer? You must always verify:
\[
g \sim g' \implies \phi(g) = \phi(g').
\]
If this check fails, your formula is not a function at all; it is an ambiguous assignment.
\end{trapbox}

Figure~\ref{32:fig-welldefinedness} illustrates this foundational check: every element within the same equivalence class must map to the exact same target value for the induced quotient map to be single-valued.

\begin{figure}[htbp]
\centering
\begin{tikzpicture}[scale=0.95, >=Latex]
  % Left Panel: Well-Defined
  \begin{scope}[xshift=0cm]
    \node[anchor=south, font=\small\bfseries, text=targetgreen!90!black] at (2.7, 2.1) {Sound: Well-Defined Map};
    \draw[rounded corners=4pt, fill=gray!3, draw=gray!25] (-0.2, -1.9) rectangle (5.6, 2.1);
    
    % Equivalence class [g]
    \draw[thick, draw=sourceblue!80!black, fill=sourceblue!12] (1.1, 0.05) ellipse (0.75 and 1.15);
    \node[font=\footnotesize\bfseries, text=sourceblue!90!black] at (1.1, 0.9) {$[g]$};
    
    \fill[sourceblue!90!black] (1.1, 0.45) circle (2.2pt);
    \node[left=3pt, font=\scriptsize\bfseries, text=sourceblue!90!black] at (1.1, 0.45) {$g$};
    
    \fill[sourceblue!90!black] (1.1, -0.35) circle (2.2pt);
    \node[left=3pt, font=\scriptsize\bfseries, text=sourceblue!90!black] at (1.1, -0.35) {$g'$};
    
    % Codomain point
    \fill[targetgreen!90!black] (4.5, 0.05) circle (2.5pt);
    \node[right=3pt, font=\footnotesize\bfseries, text=targetgreen!90!black] at (4.5, 0.05) {$y$};
    
    % Arrows
    \draw[->, thick, draw=sourceblue!70!black] (1.2, 0.45) to[bend left=12] (4.3, 0.15);
    \draw[->, thick, draw=sourceblue!70!black] (1.2, -0.35) to[bend right=12] (4.3, -0.05);
    
    \node[font=\scriptsize, align=center, text=gray!80!black] at (2.7, -1.45) {Condition: $g \sim g' \implies \phi(g) = \phi(g')$\\Output $\bar{\phi}([g]) = y$ is single-valued};
  \end{scope}

  % Right Panel: Ill-Defined
  \begin{scope}[xshift=6.6cm]
    \node[anchor=south, font=\small\bfseries, text=trapred!90!black] at (2.7, 2.1) {Flawed: Ambiguous Assignment};
    \draw[rounded corners=4pt, fill=gray!3, draw=gray!25] (-0.2, -1.9) rectangle (5.6, 2.1);
    
    % Equivalence class [g]
    \draw[thick, draw=sourceblue!80!black, fill=sourceblue!12] (1.1, 0.05) ellipse (0.75 and 1.15);
    \node[font=\footnotesize\bfseries, text=sourceblue!90!black] at (1.1, 0.9) {$[g]$};
    
    \fill[sourceblue!90!black] (1.1, 0.45) circle (2.2pt);
    \node[left=3pt, font=\scriptsize\bfseries, text=sourceblue!90!black] at (1.1, 0.45) {$g$};
    
    \fill[sourceblue!90!black] (1.1, -0.35) circle (2.2pt);
    \node[left=3pt, font=\scriptsize\bfseries, text=sourceblue!90!black] at (1.1, -0.35) {$g'$};
    
    % Codomain points
    \fill[trapred!90!black] (4.3, 0.6) circle (2.2pt);
    \node[right=3pt, font=\scriptsize\bfseries, text=trapred!90!black] at (4.3, 0.6) {$y_1 = \phi(g)$};
    
    \fill[trapred!90!black] (4.3, -0.5) circle (2.2pt);
    \node[right=3pt, font=\scriptsize\bfseries, text=trapred!90!black] at (4.3, -0.5) {$y_2 = \phi(g')$};
    
    % Arrows diverging
    \draw[->, thick, draw=trapred!80!black, dashed] (1.2, 0.45) to[bend left=12] (4.1, 0.6);
    \draw[->, thick, draw=trapred!80!black, dashed] (1.2, -0.35) to[bend right=12] (4.1, -0.5);
    
    \node[font=\scriptsize, align=center, text=trapred!90!black] at (2.7, -1.45) {Conflicting outputs: $\phi(g)$ vs $\phi(g')$\\Input $[g]$ branches to two targets!};
  \end{scope}
\end{tikzpicture}
\caption{The well-definedness criterion for quotient maps. \textbf{Left}: If every representative in the class $[g]$ maps to the same target $y$, the induced map $\bar{\phi}([g]) = y$ is a legitimate single-valued function. \textbf{Right}: If different representatives $g$ and $g'$ produce conflicting outputs $y_1 \neq y_2$, the rule assigns multiple values to the same input, violating the definition of a function.}
\label{32:fig-welldefinedness}
\end{figure}

\subsection{What Morphisms Achieve---and Their Limits}

Viewing functions as structure-preserving morphisms transformed mathematics:
\begin{enumerate}
  \item \textbf{Universal Definitions}: Concepts like products ($A \times B$), quotients ($A/\sim$), and free groups (collections of words formed by letters with no relations other than inverses canceling) are defined not by ad-hoc recipes, but by their \emph{universal mapping property}---defining an object by how every other object must map into it or out of it, guaranteeing existence and uniqueness up to unique isomorphism.
  \item \textbf{Bridges Across Disciplines}: Functors translate geometric shapes into algebraic equations. For example, algebraic topologists assign to any space its fundamental group (measuring how loops wind around holes) and homology groups (counting higher-dimensional cavities).
  \item \textbf{Invariant Decomposition}: Any complicated morphism cleanly factors into a canonical projection followed by an isomorphism and an inclusion ($f = \iota \circ \bar{f} \circ \pi$).
\end{enumerate}

\paragraph{Scope and Limitations.}
Yet arrow abstraction has genuine boundaries:
\begin{itemize}
  \item \textbf{Foundational Size Paradoxes}: In naive category theory, one speaks easily of ``the category of all sets'' or ``the category of all groups.'' But just as there is no ``set of all sets'' without triggering Russell's paradox, a collection of all categories contains itself and risks logical contradiction unless managed with larger universes or set-class distinctions (such as Gödel--Bernays set theory).
  \item \textbf{Loss of Computational Specifics}: Category theory proves that a solution exists and is unique up to unique isomorphism, but it rarely produces an explicit numerical formula or an algorithm to solve a differential equation. When an engineer needs the vibration frequencies of a bridge, categorical arrows must eventually yield to concrete matrix computations.
\end{itemize}

\paragraph{Where We Go Next in the Book.}
In this chapter, we saw that every function factors by gathering inputs with identical outputs into fibers, collapsing those fibers into a quotient. In Chapter~33 (\emph{Second Pass: Relations, Classification, and Quotients}), we inspect the exact machinery behind this bundling: the equivalence relation. We will study the three foundational axioms (reflexivity, symmetry, transitivity) that guarantee equivalence classes tile a space without gaps, and watch how quotients construct circles, cylinders, and modular arithmetic from ordinary Euclidean lines.

\subsection{Exercises}

\begin{enumerate}
  \item \textbf{Mental Check (Structure Preservation)}: For each function below, determine whether it preserves the specified group operation without writing a full proof:
  \begin{enumerate}
    \item $f: (\R, +) \to (\R, +)$ defined by $f(x) = 3x$.
    \item $g: (\R, +) \to (\R, +)$ defined by $g(x) = x + 3$.
    \item $h: (\R_{>0}, \times) \to (\R_{>0}, \times)$ defined by $h(x) = x^3$.
  \end{enumerate}
  \item For each function below, identify the fibers (the sets of inputs mapped to the same output) and describe what geometric information is forgotten:
  \begin{enumerate}
    \item $f: \R^2 \to \R$ defined by $f(x, y) = x^2 + y^2$.
    \item $g: \R^3 \to \R^2$ defined by $g(x, y, z) = (x, y)$.
    \item $\operatorname{det}: \operatorname{GL}_2(\R) \to \R^\times$, mapping invertible $2 \times 2$ matrices to their nonzero determinant.
  \end{enumerate}
  \item Let $\phi: \Z \to \Z_6$ send an integer $n$ to $n \bmod 6$. Show that $\phi$ is a ring homomorphism (preserving both addition and multiplication). What is $\ker(\phi)$? What does $\Z/\ker(\phi)$ look like?
  \item Suppose $f: A \to B$ and $g: B \to C$ are functions with composite $h = g \circ f$.
  \begin{enumerate}
    \item If $h$ is injective, prove that $f$ must be injective. Give a concrete counterexample showing $g$ need not be injective.
    \item If $h$ is surjective, prove that $g$ must be surjective. Give a counterexample showing $f$ need not be surjective.
  \end{enumerate}
  \item Consider the wrapping map $\phi: (\R, +) \to (S^1, \cdot)$ given by $\phi(t) = e^{2\pi i t}$, where $S^1 = \{z \in \C : |z| = 1\}$.
  \begin{enumerate}
    \item Prove that $\phi$ is a group homomorphism.
    \item Prove that $\ker(\phi) = \Z$.
    \item Apply the First Isomorphism Theorem to deduce that $\R/\Z \cong S^1$. Describe the geometric meaning of the cosets $t + \Z$.
  \end{enumerate}
  \item Let $\phi: (\Z_3, +) \to (\Z_5, +)$ be a group homomorphism.
  \begin{enumerate}
    \item Show that for any $x \in \Z_3$, the element $\phi(x)$ must satisfy $3 \cdot \phi(x) = 0$ in $\Z_5$.
    \item Using the fact that $\gcd(3, 5) = 1$, prove that $\phi(x) = 0$ for all $x \in \Z_3$.
    \item Conclude that the only homomorphism from $\Z_3$ to $\Z_5$ is the trivial map.
  \end{enumerate}
\end{enumerate}

\newpage
